%%=============================================================
%%  Master-Thesis status report — meeting handout
%%  Author : Esra D.
%%  Date   : 22 May 2026
%%  Style  : IEEE conference (final paper format)
%%=============================================================
\documentclass[conference]{IEEEtran}

\IEEEoverridecommandlockouts

\usepackage{cite}
\usepackage{amsmath,amssymb,amsfonts}
\usepackage{algorithmic}
\usepackage{graphicx}
\usepackage{textcomp}
\usepackage{xcolor}
\usepackage{booktabs}
\usepackage{array}
\usepackage{hyperref}
\usepackage{listings}
\usepackage{url}

\hypersetup{colorlinks=true,linkcolor=blue,citecolor=blue,urlcolor=blue}

\lstset{
  basicstyle=\ttfamily\footnotesize,
  breaklines=true,
  columns=fullflexible,
  frame=single,
  framesep=3pt,
  xleftmargin=4pt,
  keepspaces=true
}

\def\BibTeX{{\rm B\kern-.05em{\sc i\kern-.025em b}\kern-.08em
    T\kern-.1667em\lower.7ex\hbox{E}\kern-.125emX}}

\begin{document}

\title{Stereotype-Guided Q-Learning for Smart-Building \\Illuminance Control: \\Thesis Status Report
\\ \large (Working document --- supervisor meeting, 22~May~2026)}

\author{\IEEEauthorblockN{Esra D.}
\IEEEauthorblockA{M.Sc.\ Thesis project --- MT-Esra\\
JaCaMo-based Multi-Agent System\\
Submission deadline: 30 June 2026}}

\maketitle

\begin{abstract}
This report summarises the current state of my Master's thesis project
\emph{MT-Esra}, a JaCaMo multi-agent system that investigates whether
\emph{stereotype-guided Q-learning initialisation}---using BRICK + Elementary
Ontology knowledge encoded as soft Q-table priors---yields measurably better
control of a simulated multi-zone lighting environment than (i)~a purely
ontology-driven rule-based agent and (ii)~a standard tabular Q-learner with
zero-initialised Q-tables. The report covers the project's architecture
(agents, environment artifacts, ontology), the experimental protocol over
seven simulator profiles (one ``clean'' lab and six labs that each inject a
single systematic ontology gap, $W1\ldots W6$), the benchmark results
obtained so far, the conclusions drawn from those results, the three research
extensions I have implemented in the last week
(multi-seed pipeline with bootstrap confidence intervals,
Potential-Based Reward Shaping, and adaptive trust calibration of the
stereotype priors), and a six-week plan that ends with the 10-page
IEEE-format conference paper which is the contractual final deliverable.
\end{abstract}

\begin{IEEEkeywords}
Multi-agent systems, JaCaMo, reinforcement learning, Q-learning,
stereotypes, knowledge graphs, BRICK, smart buildings, Web of Things.
\end{IEEEkeywords}

%=============================================================
\section{Project Overview and Motivation}
\label{sec:overview}

\subsection{Research Question}
Pure reinforcement learning agents need many interaction samples to converge
because they start with no prior over the environment. Pure rule-based
agents start with strong domain knowledge but cannot recover when that
knowledge is incomplete or wrong. The thesis asks one question:

\begin{quote}
\emph{Given an explicit, machine-readable building ontology
(BRICK + Elementary), can we use it to pre-shape a Q-learner's value
function so that the agent converges faster and acts more safely
\textbf{without} losing the ability to recover when the ontology is
incomplete or wrong?}
\end{quote}

The contribution we target is a \emph{quantitative} answer: under which
ontology-gap conditions stereotype-guided initialisation helps, hurts, or
is neutral, and a small set of mechanisms (PBRS, adaptive trust) that
\emph{bound} the worst-case damage when the ontology is wrong.

\subsection{Why JaCaMo / multi-agent}
The lighting controller is naturally a hybrid agent: it must reason
symbolically about zones, actuators and topology (a job for AgentSpeak +
SPARQL), and it must learn quantitative trade-offs from data (a job for a
Q-learning artifact). JaCaMo's clean separation between Jason
\emph{agents} and CArtAgO \emph{artifacts} lets us put both in one program
without forcing either side to compromise: the agent code stays declarative
and ontology-driven, while the learning logic lives in plain Java where
linear algebra is convenient.

\subsection{Why a smart-building lighting testbed}
Smart buildings are an ideal stress test for ``ontology helps RL'' claims:
\begin{itemize}
  \item The domain has a widely-adopted ontology (BRICK) plus a process
        ontology (Elementary) that can express \emph{stereotypes}
        of components (e.g.\ ``a blind \emph{mediates} sunshine onto a
        zone'') \cite{cecconi2023}.
  \item State and action spaces are small enough for tabular Q-learning
        to be a sound baseline (no neural-network confounders).
  \item Real failure modes (hidden cross-zone bleed, power caps,
        thermal coupling, dynamic topology) are easy to model in
        Node-RED and map naturally onto six well-defined ontology
        gaps $W1\ldots W6$.
\end{itemize}

%=============================================================
\section{System Architecture}
\label{sec:arch}

Figure~\ref{fig:arch} (drawn during the meeting) shows the runtime stack.
At a glance:

\begin{verbatim}
+-----------------------------------------+
|  Jason agents (AgentSpeak)              |
|  - illuminance_controller_agent.asl     |  rule-based
|  - illuminance_controller_agent_ql.asl  |  Q-learning (train)
|  - illuminance_controller_agent_bench   |  benchmark driver
|  - lab_profiles.asl                     |  profile registry
+--------------------+--------------------+
                     |  CArtAgO operations
+--------------------v--------------------+
|  Environment artifacts (Java)           |
|  LabEnvironment   - WoT/HTTP I/O        |
|  OntologyArtifact - SPARQL on KG        |
|  QLearner         - tabular Q-learning  |
|  StereotypeReasoner - prior penalties   |
|  StereotypeLearner  - online effects    |
|  BenchmarkLogger  - per-step metrics    |
+--------------------+--------------------+
                     |  HTTP (Thing Description)
+--------------------v--------------------+
|  Node-RED simulator (ports 1881..1888)  |
|  one flow per lab profile               |
+-----------------------------------------+
\end{verbatim}

\subsection{Agents (Jason)}
\begin{itemize}
  \item \textbf{Rule-based agent} --- discovers zones and actuators
        dynamically via SPARQL over the BRICK + Elementary KG; no
        component name is hard-coded. It is the symbolic baseline.
  \item \textbf{Q-learning agent} --- the same skeleton as the
        rule-based agent for state reading and action dispatch, but
        delegates action choice to \texttt{QLearner.getActionFromState}.
        A single belief, \verb|use_stereotypes(true|false)|, toggles
        prior-guided initialisation.
  \item \textbf{Benchmark agent} --- runs a fixed scenario set, in one
        of three modes (\verb|rule_based|, \verb|ql_false|, \verb|ql_true|),
        and per step calls \verb|BenchmarkLogger.recordStepDetailRich(...)|
        which also runs the $W1\ldots W6$ weakness classifier against
        the ontology's prediction.
  \item \textbf{Lab profile registry} --- \verb|lab_profiles.asl| is the
        single source of truth: TD path, ontology files, scenarios path,
        simulator port, illuminance discretisation bounds, zone targets,
        sunshine probability, weakness flags, Q-table suffix, training
        budget. Switching labs is one line.
\end{itemize}

\subsection{Environment Artifacts (Java/CArtAgO)}
\begin{itemize}
  \item \texttt{LabEnvironment} --- parses a W3C Web-of-Things Thing
        Description (TTL) and dispatches HTTP actions to Node-RED.
        State is read via REST.
  \item \texttt{OntologyArtifact} --- Apache Jena SPARQL endpoint over
        \verb|lab-ontology.ttl| + \verb|wot-mappings.ttl|. Provides
        dynamic zone discovery, quantity discovery and ``best-action''
        queries used by the rule-based agent.
  \item \texttt{StereotypeReasoner} --- plain Java; computes
        \emph{initial Q-table penalties} per (state, action, zone)
        from Cecconi et al.'s component stereotypes
        \cite{cecconi2023}. Hard penalty $-100$ for redundant actions,
        soft penalty $-100\cdot(1-p_{\text{sat}})$ for IV (instrumental
        variable) mismatch, plus cross-zone overshoot and shared-actuator
        conflict penalties.
  \item \texttt{QLearner} --- decomposed tabular Q-learning
        \(Q_z[s][a]\), one table per zone. Update is standard
        $\epsilon$-greedy with $\alpha=0.1,\gamma=0.9$, $\epsilon$ decays
        $0.3 \to 0.01$. \emph{Per-zone reward} is computed inside
        \verb|computeZoneReward| (per-step cost $-1$, gap-progress
        $\times 40$, terminal bonus $+200$, regression penalty $-200$,
        no-effect $-10$, stagnation $-5$). Reward is clipped at
        $\pm 50$ since 2026-05-21 (Phase~A fix).
  \item \texttt{StereotypeLearner} --- observes
        $(s,a,s')$ triples online and uses Welford's algorithm to keep
        per-effect running mean and variance. Effects above a
        Z-score threshold are emitted as Turtle to
        \verb|learned_stereotypes_*.ttl| (gap discovery output).
  \item \texttt{BenchmarkLogger} --- per-scenario summary
        \verb|benchmark_results_<mode>.csv| (goal rate, steps,
        deviation, energy, $W1$--$W6$ counters) and per-step
        \verb|bench_step_log_<mode>.csv|. Phase~B added the
        \verb|StuckFired| column for the anti-stuck guard.
\end{itemize}

\subsection{State / action spaces}
The state vector is read directly from \verb|wot-mappings.ttl|; for the
four-zone lab it has 13 slots
$\bigl[\,Z_{1..4}\text{Level},
       Z_{1..4}\text{Light},
       Z_{1..4}\text{Blinds},
       \text{Spotlight},
       \text{Sunshine}\,\bigr]$.
Continuous lux is discretised to four ranks ($0\!=\!<\!75$,
$1\!=\!75\text{--}200$, $2\!=\!200\text{--}400$, $3\!=\!\geq 400$).
There are 21 actions: \verb|DO_NOTHING|, \verb|ON/OFF| per
\verb|Z<n>Light|, and \verb|UP/DOWN| per \verb|Z<n>Blinds|.

%=============================================================
\section{Experimental Protocol}
\label{sec:protocol}

\subsection{Lab profiles}

\begin{table}[h]
\caption{Simulator profiles. Each ``customN'' lab injects exactly one
systematic ontology gap; \texttt{custom8} is the clean lab where the
ontology is correct.}
\label{tab:profiles}
\centering
\footnotesize
\begin{tabular}{@{}llp{4.5cm}@{}}
\toprule
Profile & Gap & Description \\
\midrule
\texttt{custom8} & ---  & Clean 4-zone lab; ontology is correct.\\
\texttt{custom2} & $W1$ & Missing-stereotype: corridor light bleeds $+30$ lux to all zones (unmodelled).\\
\texttt{custom3} & $W2$ & Inversion: \texttt{Blind\_Z2} effect flips sign when \texttt{Sunshine}~$\geq 500$ lux.\\
\texttt{custom4} & $W3$ & Dynamics: per-lamp ramp + hysteresis (no immediate effect).\\
\texttt{custom5} & $W4$ & Shared resource: 7-unit power cap silently drops lowest-priority lamp.\\
\texttt{custom6} & $W5$ & Dynamic topology: \texttt{Partition23} toggles every 5 ticks.\\
\texttt{custom7} & $W6$ & Multi-objective heat: every ON lamp raises zone temperature by $+0.1\,^{\circ}\!$C per tick.\\
\bottomrule
\end{tabular}
\end{table}

\subsection{Modes under test}
For each profile we run three agent modes:
\begin{itemize}
  \item \verb|rule_based| --- pure symbolic baseline,
  \item \verb|ql_false| --- standard Q-learning, zero-init,
  \item \verb|ql_true|  --- stereotype-guided Q-learning.
\end{itemize}
The QL modes are first trained
(\verb|num_episodes=2500| for the ``paper'' profile, $50$ for ``dev'')
against a separate \verb|train_scenarios_<profile>.json| set,
then benchmarked against \emph{held-out} \verb|scenarios_<profile>.json|
(28 scenarios) with $\epsilon=0$ (purely greedy).

\subsection{Pipeline orchestration}
The full pipeline is driven by two PowerShell scripts:
\begin{itemize}
  \item \verb|run_full_project.ps1| --- sequential reference driver
        (start simulators, preflight, train, benchmark, analysis).
  \item \verb|run_full_project_parallel.ps1| --- worktree-clone
        based parallel driver. Each profile/stereo-mode pair is
        trained in its own clone via
        \verb|Start-Process -EncodedCommand|; each profile/bench-mode
        pair is benchmarked via \verb|Start-Job|.
\end{itemize}
The runtime configuration lives in \verb|config/run_config.json|
(profiles \verb|dev|/\verb|paper|, port map, suffix map, learning block).

%=============================================================
\section{Results So Far}
\label{sec:results}

The current full benchmark sweep over \verb|custom2|--\verb|custom8|
ran on 21~May 2026 (archive
\verb|benchmark/results.archive_2026-05-21_2207_v1|). The numbers below are
from the clean lab \verb|custom8| (28 scenarios $\times$ 2 runs $= 56$
benchmark episodes), where the ontology is correct and stereotypes
\emph{should} therefore help.

\begin{table}[h]
\caption{Benchmark on clean lab \texttt{custom8} (current state).
Goal rate is the fraction of scenarios in which all zone targets were
reached within the 20-step budget. Lower is better for the other columns.}
\label{tab:custom8}
\centering
\footnotesize
\begin{tabular}{@{}lrrrr@{}}
\toprule
Mode & Goal rate & Avg steps & Avg deviation & Avg energy\\
\midrule
\texttt{rule\_based} & \textbf{0.464} & 14.5 & 42.8  & 50.5 \\
\texttt{ql\_false}   & 0.071          & 19.2 & 56.3  & 64.4 \\
\texttt{ql\_true}    & 0.018          & 19.8 & 99.1  & 28.5 \\
\bottomrule
\end{tabular}
\end{table}

\subsection{What the results say honestly}

Three observations dominate:

\paragraph{1. The symbolic baseline is still the strongest single agent
on the clean lab.} \texttt{rule\_based} reaches the goal in 46\% of
scenarios in 14.5 steps; the Q-learners reach it in 7\% / 2\% in 19+
steps. The cause is structural, not a tuning issue: a 20-step horizon
on a 13-dimensional state space with 21 actions is genuinely small for
tabular Q-learning to converge to a near-optimal greedy policy in
2500 training episodes when the training distribution does not perfectly
cover the test distribution.

\paragraph{2. \texttt{ql\_true} currently loses to \texttt{ql\_false} on
the clean lab.} This is the inversion of the hypothesis and is the
single most important open issue. Diagnosis: the soft priors are
\emph{always} applied during greedy selection (they do not decay to
zero), and on the clean lab the prior's sunshine-conditioned penalty on
blind actions appears to discourage exploration of action combinations
that the zero-init Q-learner does try. The Phase~A fix (reward clipping
to $\pm 50$, prior decay floor $0.1$, prior scale $\times 5$, seeded
RNG XOR) corrected reward saturation but did not, by itself, recover
the expected ordering.

\paragraph{3. Energy efficiency does break for \texttt{ql\_true}.}
\texttt{ql\_true} spends \emph{half} the energy of \texttt{ql\_false}
(28.5 vs.\ 64.4) and a comparable fraction of \texttt{rule\_based}.
Even with a lower goal rate, the stereotype prior is clearly
suppressing the ``turn everything on'' attractor of the zero-init
Q-learner. Energy-efficient failure is a meaningful intermediate
result for the paper.

\subsection{Weakness fires}
The benchmark CSV also reports the $W1\ldots W6$ classifier counts.
On \verb|custom8| (clean lab) we still see large \verb|W3| (delayed
effect) counts in both QL modes because the rank-discretisation step
sometimes lags one tick behind the actuator flip; this is a known
artefact of the discrete state representation and is not informative.
The interesting weakness columns ($W1,W2,W4,W5,W6$) are zero on the
clean lab, which is exactly what the classifier should report.

%=============================================================
\section{Interpretation and Research Findings}
\label{sec:findings}

What can already be claimed for the paper:

\begin{enumerate}
  \item \textbf{An ontology-driven symbolic baseline is competitive
        with tabular Q-learning on small-horizon multi-zone lighting
        control.} Across all seven labs, \texttt{rule\_based} is
        within range of (and on \texttt{custom8} clearly above) both
        Q-learning variants on goal rate. This is itself a publishable
        observation, because it counters the implicit assumption in
        much of the smart-building RL literature that learning is
        automatically better than rule-based control.

  \item \textbf{Stereotype priors materially change the action
        distribution.} \texttt{ql\_true} uses roughly half the energy
        of \texttt{ql\_false} on \texttt{custom8} and the same trend
        holds on the weakness labs. The prior is doing its job
        suppressing useless actions; the issue is that it is also
        suppressing exploration of useful action combinations on the
        clean lab.

  \item \textbf{The Cecconi et al.\ stereotype formalism is
        operationalisable end-to-end.} I have shown that the
        \emph{Mediates}, \emph{Causes}, \emph{IV} primitives can be
        compiled from a BRICK + Elementary KG into per-zone Q-table
        initial values via SPARQL, used at runtime to mask actions,
        and updated online via \texttt{StereotypeLearner}. That
        pipeline did not exist in the literature before this thesis.

  \item \textbf{The six weakness labs are a reusable
        ontology-stress benchmark.} Each lab pins exactly one of
        $W1\ldots W6$ and the simulator flows are reproducible; the
        thesis ships them as a benchmark suite.
\end{enumerate}

What is still open:

\begin{itemize}
  \item The expected ordering \verb|ql_true| $>$ \verb|ql_false|
        on the clean lab does not hold yet in numbers. The
        extensions described in
        Section~\ref{sec:extensions} are designed to fix this.
  \item No confidence intervals: a single-seed sweep cannot
        distinguish ``\texttt{ql\_true} is genuinely worse'' from
        ``the seed was unlucky''. Section~\ref{sec:extensions} fixes
        this too.
\end{itemize}

%=============================================================
\section{Implemented Research Extensions}
\label{sec:extensions}

In the last week I implemented three extensions, all opt-in and all
preserving the existing default behaviour. The Java code builds clean
(\verb|./gradlew compileJava|); the Python and PowerShell parts
parse-check clean.

\subsection{(E1) Multi-seed pipeline with bootstrap CIs}
\label{sec:ext-seeds}

\paragraph{Motivation.} A single training run is one sample from a
high-variance distribution. Any claim like ``\texttt{ql\_true} is
better than \texttt{ql\_false}'' is meaningless without an interval.

\paragraph{Implementation.}
\begin{itemize}
  \item \verb|QLearner.java| has a new \verb|RUN_SEED| system property
        (\verb|-Prun.seed=N|). At configuration time
        $$
        \texttt{baseSeed} \leftarrow \texttt{baseSeed}\;\oplus\;\text{mix64}(N),
        $$
        where mix64 is the SplitMix64 finalizer (so small $N$ values
        spread across all 64 bits).
  \item \verb|run_full_project.ps1| has a new \verb|-RunSeed| parameter
        that forwards as \verb|-Prun.seed=N|.
  \item \verb|run_full_project_parallel.ps1| has a new \verb|-Seeds 1,2,3,4,5|
        parameter. Phases D--G (train + benchmark + per-clone cleanup)
        are wrapped in a per-seed loop; each seed's
        \verb|benchmark/results| is archived to
        \verb|benchmark/results_seed<N>/| and the clone is wiped clean
        between seeds.
  \item \verb|analysis/sweep_report.py| has \verb|--seeds-mode| that
        detects \verb|results_seed*/| siblings, computes
        non-parametric 95\% bootstrap CIs (numpy with
        \verb|default_rng(0xC1)|, stdlib fallback) over
        \verb|goal_rate|, \verb|avg_steps|, \verb|avg_dev|,
        \verb|avg_energy|, \verb|avg_wasted|, and writes
        \verb|summary_table_ci.csv| with mean and CI bounds per
        (profile, mode) cell.
\end{itemize}
\paragraph{Result.} Every numerical claim in the paper will be
reported as \emph{mean $\pm$ 95\% CI over 5 seeds}.

\subsection{(E2) Potential-Based Reward Shaping (PBRS)}
\label{sec:ext-pbrs}

\paragraph{Motivation.} The sparse $+200$ terminal reward on a 20-step
horizon means a single transition with a non-zero gap delta is the
only training signal in 90\% of steps. PBRS adds a dense \emph{shaping}
signal that is provably policy-invariant
\cite{ng1999policy}: it changes how fast the agent learns, not what
the optimal policy is.

\paragraph{Implementation.}
A new system property \verb|reward.shaping=pbrs| activates the
shaping term inside the per-zone loop of
\verb|QLearner.calculateQ|:
\begin{eqnarray*}
\Phi_z(s)  &=& -|\text{level}_z(s) - \text{target}_z|,\\
F(s,a,s')  &=& \gamma\,\Phi_z(s') - \Phi_z(s),\\
r'         &=& r_{\text{clipped}} + F(s,a,s')/N_{\text{zones}}.
\end{eqnarray*}
The factor $1/N_{\text{zones}}$ keeps the per-zone update on the same
scale as the clipped per-zone reward; the policy-invariance proof of
Ng--Harada--Russell still holds for any positive scaling.

\paragraph{Expected effect.} Faster convergence on \emph{all} labs;
no change to the policy that an infinitely-trained agent would
converge to. This is the ``smallest code change with the strongest
theoretical guarantee'' --- the change is roughly 20 lines.

\subsection{(E3) Adaptive trust calibration of stereotype priors}
\label{sec:ext-trust}

\paragraph{Motivation.} The fundamental honest objection to E1+E2 is:
\emph{``what if the ontology is wrong?''} The current Phase~A code
uses a fixed linear decay (\verb|stereo.priorDecayEpisodes=10000|
with floor 0.1) that applies the same trust schedule whether the
ontology is correct (\texttt{custom8}) or wrong ($W1\ldots W6$).
That is unsatisfying because it cannot exploit the information that
\texttt{custom8} \emph{should} keep trusting the prior.

\paragraph{Implementation.}
Activated by \verb|stereo.adaptiveTrust=true|. For every observed
transition I ask the \texttt{StereotypeReasoner}:
\emph{``would you have predicted the sign of the observed zone-level
$\Delta$ correctly?''} I maintain a per-action running calibration
$$
c_a = \frac{\#\text{slots where pred.\ sign matched observed sign}}
            {\#\text{slots where pred.\ sign was non-zero}},
$$
stored as cumulative \verb|actionCalSum[a]| and \verb|actionCalN[a]|.
After at least \verb|stereo.adaptiveTrust.minSamples| (default 50)
observations, the prior used in \verb|greedyAction| becomes
$$
\text{effectivePrior}(a) =
   \text{priorWeight}\cdot p_a \cdot
   \max\bigl(\text{floor},\;c_a\bigr),
$$
where the floor (default 0.1) prevents complete collapse and the
multiplier is 1.0 when there are not yet 50 observations.

\paragraph{Expected effect.}
On the clean lab (\texttt{custom8}), $c_a \to 1$ and the prior keeps
its full weight throughout training. On a weakness lab (e.g.\
\texttt{custom2}), the $W1$-affected actions accumulate
\emph{prediction mismatch} and $c_a$ collapses toward the floor,
recovering the \texttt{ql\_false} baseline. This addresses the
\emph{exact} symmetric failure mode of the current code.

\subsection{Why these three together}
E1 (seeds) gives every claim a CI; E2 (PBRS) gives faster, provably
policy-equivalent learning; E3 (adaptive trust) gives a graceful
fallback when the ontology is wrong. Together they let the paper make
a falsifiable, two-sided claim: \emph{stereotypes never hurt and help
when correct}, with $95$\% CIs.

%=============================================================
\section{Plan for the Remaining Six Weeks}
\label{sec:plan}

\subsection{Week 1 (23--29 May): validate the extensions empirically}

\begin{enumerate}
  \item Smoke test the multi-seed driver:
        \begin{lstlisting}[language=PowerShell]
.\run_full_project_parallel.ps1 `
   -RunMode dev -Seeds 1,2 -MaxParallel 2 -Fresh
        \end{lstlisting}
        Confirm \verb|benchmark/results_seed1|,
        \verb|results_seed2|, and \verb|analysis/out/summary_table_ci.csv|.
  \item Single-profile \verb|custom8| dev sweep with PBRS on, baseline
        off; check the learning curve shifts left (faster
        convergence).
  \item Single-profile \verb|custom2| dev sweep with adaptive trust
        on; check that \verb|ql_true| recovers toward \verb|ql_false|
        on the $W1$-affected actions.
\end{enumerate}

\subsection{Week 2 (30~May--5~Jun): paper-quality multi-seed sweep}
Launch the 5-seed paper-mode sweep over all seven profiles:
\begin{lstlisting}[language=PowerShell]
.\run_full_project_parallel.ps1 `
   -RunMode paper -Seeds 1,2,3,4,5 -MaxParallel 3
\end{lstlisting}
Estimated wall-clock: 5 seeds $\times \sim 50$~h $/$ 3 parallel
$\approx$ 4 calendar days. The output is the headline number for the
paper.

\subsection{Week 3 (6--12~Jun): ablations}
Three controlled ablations, each one extension turned off in turn:
\begin{itemize}
  \item \verb|reward.shaping=none| with E1+E3 (does PBRS matter?),
  \item \verb|stereo.adaptiveTrust=false| with E1+E2 (does adaptive
        trust matter on \texttt{custom2}--\texttt{custom7}?),
  \item \verb|stereo.priorScale=0.0| with E1+E2+E3 (sanity check:
        with no prior at all, does \texttt{ql\_true} collapse onto
        \texttt{ql\_false}?).
\end{itemize}

\subsection{Week 4 (13--19~Jun): write paper sections 1--4}
Sections \emph{Introduction}, \emph{Background and Related Work}
(Cecconi 2023, BRICK, JaCaMo, classical RL shaping references),
\emph{System}, \emph{Experimental Protocol}. Reuse text from this
report and from \verb|README.md|. Target $\approx 5$ pages of the
$10$-page budget.

\subsection{Week 5 (20--26~Jun): write paper sections 5--8 + figures}
\emph{Results}, \emph{Discussion}, \emph{Limitations and Threats to
Validity}, \emph{Conclusion}. All figures regenerated from the
seeds-mode output: bootstrap-CI error bars on every bar chart,
per-weakness heatmap, learning curves with confidence bands.
Target the remaining $\approx 5$ pages.

\subsection{Week 6 (27--30~Jun): polish, anonymous review pass, submit}
Internal review pass with supervisor, IEEE checklist
(PDF/A, copyright, references, figure resolution), final submission
on 30~June.

\subsection{Risk register and mitigations}
\begin{itemize}
  \item \emph{Risk:} the paper-mode 5-seed sweep does not finish in
        Week~2. \emph{Mitigation:} reduce to 3 seeds (still
        statistically usable) or to 4 weakness labs (\texttt{custom2},
        \texttt{custom4}, \texttt{custom5}, \texttt{custom7} + clean).
  \item \emph{Risk:} adaptive trust fails to recover \texttt{ql\_false}
        baseline on weakness labs. \emph{Mitigation:} the floor and
        minSamples are config knobs; if neither value works, fall back
        to the simpler linear schedule already in place and report
        that result honestly.
  \item \emph{Risk:} parallel driver deadlocks under contention.
        \emph{Mitigation:} the sequential
        \verb|run_full_project.ps1| is a drop-in fallback;
        slower but verified.
\end{itemize}

%=============================================================
\section{Discussion Points for the Meeting}
\label{sec:meeting}

I would like the supervisor's input on three questions:

\begin{enumerate}
  \item \textbf{Honesty about the current ordering.} The paper will
        report that \texttt{ql\_true} currently loses to
        \texttt{ql\_false} on the clean lab and that E1--E3 are
        designed to fix that. Should we frame the paper as
        ``stereotypes help when calibrated'' (positive framing) or
        as ``a recipe for safe use of ontology priors in RL''
        (defensive framing)?
  \item \textbf{Scope of the weakness benchmark.} Six gap labs may be
        too much for 10 pages. Acceptable to reduce to four
        representative labs in the paper while keeping all seven in
        the thesis appendix?
  \item \textbf{Comparison to non-tabular baselines.} A reviewer will
        ask why not DQN. My answer: tabular avoids
        function-approximation confounds and the state space is
        small enough that tabular is the correct method here.
        Do you agree this is enough?
\end{enumerate}

%=============================================================
\section{Conclusion}
\label{sec:conclusion}

The system, the experimental protocol, and the analysis pipeline are
fully in place. The remaining six weeks are dedicated to running the
multi-seed paper sweep, three controlled ablations, and writing the
10-page IEEE conference paper. The three extensions (E1 seeds + CIs,
E2 PBRS, E3 adaptive trust) directly target the two open issues
identified by the current single-seed sweep: lack of confidence
intervals and the inverted clean-lab ordering. All extensions are
implemented, build clean, and ship as opt-in flags in
\verb|config/run_config.json| so the existing experiments remain
exactly reproducible.

%=============================================================
\begin{thebibliography}{9}

\bibitem{cecconi2023}
F.~Cecconi \emph{et al.}, ``Reasoning about Physical Processes in
Buildings through Component Stereotypes,'' in \emph{Proc.\ CPS-IoT
Week 2023}, 2023.

\bibitem{ng1999policy}
A.~Y.~Ng, D.~Harada, and S.~Russell, ``Policy invariance under reward
transformations: Theory and application to reward shaping,'' in
\emph{Proc.\ ICML}, 1999, pp.\ 278--287.

\bibitem{brick}
B.~Balaji \emph{et al.}, ``Brick: Towards a unified metadata schema
for buildings,'' in \emph{Proc.\ ACM BuildSys}, 2016.

\bibitem{jacamo}
O.~Boissier, R.~H. Bordini, J.~F. H\"ubner, A.~Ricci, and A.~Santi,
``Multi-agent oriented programming with JaCaMo,'' \emph{Science of
Computer Programming}, vol.~78, no.~6, pp.\ 747--761, 2013.

\bibitem{wot}
W3C, ``Web of Things (WoT) Thing Description 1.1,'' W3C
Recommendation, 2023. [Online]. Available:
\url{https://www.w3.org/TR/wot-thing-description11/}

\end{thebibliography}

\end{document}
